\section{Introduction}
\label{sec:intro}

How much web text, code, encyclopedic text, and scientific writing should a pre-training corpus contain?
This choice is one of the few levers that changes what a language model can do at a fixed compute budget, yet practitioners mostly set it by heuristics.
A simple, predictable law linking domain proportions to performance would let us choose a mixture from a handful of small runs instead of trial and error at scale.

Recent work shows that such laws exist for validation loss.
\citet{ye2025datamixing} fit an exponential function of the mixture proportions to validation loss, and \citet{shukor2025scaling} extend scaling laws to include domain weights.
RegMix \citep{liu2025regmix} regresses a loss target on random mixtures and reports that web text correlates best with downstream performance, while \citet{aryabumi2025code} show that adding code to pre-training helps reasoning tasks.

However, these studies leave the \emph{shape} of the relationship between a domain's share and downstream performance open.
Loss-level laws do not tell us whether task accuracy follows the same curve, and the code result rests on a handful of mixtures, not on a dose--response curve.
Folk beliefs about mixtures are also more specific than what has been tested: that returns from a domain diminish quickly, that code drives reasoning, and that too much web text dilutes knowledge.

\para{Our main question is whether these beliefs survive a controlled test.}
We turn them into four claims.
(H1) Performance is log-linear in a domain's share.
(H2) Returns saturate at about 20\% for any single domain.
(H3) Code share is the strongest positive correlate of reasoning scores, with $r > 0.7$.
(H4) Web text above 50\% degrades knowledge-intensive tasks.

We test these claims with a dose--response study.
We train 81 GPT models with 17M parameters from scratch, each on 60M tokens drawn from a controlled mixture of four domains: web text, Wikipedia, Python code, and arXiv papers.
The design has three parts (\secref{sec:method}): single-domain sweeps that vary one share from 0 to 1, 16 random mixtures that give an unbiased sample for correlations and a held-out test set for mixture laws, and three seeds at 14 mixtures that give a noise floor.
We evaluate every model on held-out loss per domain, on zero-shot and fine-tuned knowledge question answering, and on fine-tuned reasoning probes.

The hypothesis holds for loss and essentially nowhere else.
In-domain loss is log-linear in the domain's share: the log law has 2.3 to 22 times lower leave-one-out error than a linear law, with $R^2$ between 0.98 and 0.9996.
But there is no saturation at 20\%; loss falls a further 0.36--0.55 nats from 20\% to 100\%, which is 59--98 seed standard deviations.
Downstream accuracy is almost flat.
Code share correlates with the reasoning composite at $r = -0.10$ on the random mixtures (95\% CI $[-0.55, 0.57]$), and 75\% web text is indistinguishable from 25\% or 50\% on both knowledge measures.
The largest downstream effect is in the opposite direction from H3: a pure-code corpus loses 7.4 points of fine-tuned \sciq accuracy.

These results come with a clear boundary.
Our models are four to five orders of magnitude smaller than production language models, and their zero-shot accuracy is 2--3 points above chance.
The downstream nulls therefore show that the effects are small relative to our noise at this scale, not that they are absent at large scale.

In summary, our contributions are:
\begin{itemize}[leftmargin=*,itemsep=0pt,topsep=0pt]
    \item We design a dose--response study of pre-training mixtures with 81 models on 53 mixtures, combining single-domain sweeps, random mixtures, and seed replicates so that each claim has a test and a noise estimate.
    \item We show that in-domain loss follows a log-share law with no saturation at 20\%, and that the first 2\% of a domain delivers 43--67\% of its total loss benefit.
    \item We show that a log-share law fitted on 29 mixtures predicts loss on 16 unseen mixtures with correlation 0.94--0.9996, and we compare it with the exponential law of \citet{ye2025datamixing}.
    \item We find that downstream accuracy does not follow the hypothesized laws at this scale: it is flat within seed noise across interior mixtures, and neither the code--reasoning claim nor the web--knowledge claim is supported.
\end{itemize}

We review related work in \secref{sec:related}, describe the setup in \secref{sec:method}, report results in \secref{sec:results}, and discuss interpretation and limitations in \secref{sec:discussion}.
